\begin{pastexam}[label={px:ust-107-ee-ds-08}]{臺聯大電機類 107 學年度 第八題}
（問答；10 分）

(10\%) What is the adjacency list of the following graph?

\begin{center}
\begin{tikzpicture}[v/.style={circle,fill=black!75,minimum size=7mm,inner sep=0pt,text=white,font=\small},>={Stealth[length=2.2mm]},thick]
  \node[v] (0) at (0,0) {0};
  \node[v] (1) at (-1.6,-1.7) {1};
  \node[v] (2) at (2.0,-1.7) {2};
  \node[v] (3) at (2.0,-3.6) {3};
  \node[v] (4) at (3.8,-3.6) {4};
  \node[v] (5) at (2.0,-5.5) {5};
  \draw[->] (0) to[out=-120,in=60] (1);
  \draw[->] (0) to[out=-55,in=125] (2);
  \draw[->] (1) -- (2);
  \draw[->] (2) -- (3);
  \draw[->] (2) to[out=-50,in=115] (4);
  \draw[->] (3) -- (4);
  \draw[->] (3) -- (5);
  \draw[->] (3) to[out=180,in=-60] (1);
  \draw[->] (1.west) .. controls (-2.7,-1.4) and (-2.8,1.4) .. (-1.2,0.85) .. controls (-0.7,0.7) and (-0.45,0.6) .. (0.120);
  \draw[->] (4.east) .. controls (5.1,-2.5) and (4.9,0.9) .. (2.6,0.85) .. controls (1.5,0.85) and (0.7,0.6) .. (0.60);
\end{tikzpicture}
\end{center}
\end{pastexam}

\begin{pastsoln}
先把圖上的有向邊逐條讀出來（箭頭指向的是終點）：
\[
0 \to 1,\ 0 \to 2,\quad 1 \to 0,\ 1 \to 2,\quad 2 \to 3,\ 2 \to 4,\quad 3 \to 1,\ 3 \to 4,\ 3 \to 5,\quad 4 \to 0 ,
\]
共 $10$ 條。 其中 $1 \to 0$ 是從 $1$ 左邊繞上去的弧，$4 \to 0$ 是從 $4$ 右邊繞上去的大弧，$3 \to 1$ 是從 $3$ 往左上
彎到 $1$ 的弧；$5$ 沒有出邊。

相鄰串列（資結書第~21.2~節）：每個頂點一個串列，放它的出邊指到的頂點，這裡依編號由小到大排。
\begin{center}
\begin{tikzpicture}[x=1.15cm,y=0.75cm,
  hd/.style={draw,thick,minimum width=7mm,minimum height=5mm,inner sep=1pt},
  nd/.style={draw,thick,minimum width=6mm,minimum height=5mm,inner sep=1pt},
  ar/.style={-{Stealth[length=2mm]},thick}]
\foreach \v/\lst in {0/{1,2}, 1/{0,2}, 2/{3,4}, 3/{1,4,5}, 4/{0}} {
  \node[hd] (h\v) at (0,-\v) {\v};
  \coordinate (prev) at (h\v.east);
  \foreach \w [count=\i] in \lst {
    \node[nd] (n\v\i) at (\i+0.4,-\v) {\w};
    \draw[ar] (prev) -- (n\v\i.west);
    \coordinate (prev) at (n\v\i.east);
  }
  \node[anchor=west,font=\small] at ([xshift=1mm]prev) {\texttt{NULL}};
}
\node[hd] (h5) at (0,-5) {5};
\node[anchor=west,font=\small] at ([xshift=1mm]h5.east) {\texttt{NULL}};
\end{tikzpicture}
\end{center}
寫成文字就是：
\begin{center}
\begin{tabular}{ll}
\toprule
頂點 & 相鄰串列 \\
\midrule
0 & $1 \to 2$ \\
1 & $0 \to 2$ \\
2 & $3 \to 4$ \\
3 & $1 \to 4 \to 5$ \\
4 & $0$ \\
5 & （空） \\
\bottomrule
\end{tabular}
\end{center}
每條邊在串列裡恰好出現一次，串列節點總數 $= |E| = 10$；出度依序是 $2, 2, 2, 3, 1, 0$，加起來也是 $10$。
\end{pastsoln}
